\chapter{شخصی‌سازی اعلان خط فرمان (Shell Prompt)}

در این فصل، به بررسی جزئیاتی به ظاهر ساده، یعنی اعلان خط فرمان (\en{Prompt}) خواهیم پرداخت. این واکاوی، برخی از سازوکارهای درونی پوسته و نرم‌افزار شبیه‌ساز ترمینال (\en{Terminal Emulator}) را آشکار می‌سازد. مانند بسیاری از مؤلفه‌ها در لینوکس، اعلان خط فرمان نیز به‌شدت انعطاف‌پذیر و قابل شخصی‌سازی است. اگرچه تاکنون کمتر به این موضوع پرداخته‌ایم، اما پس از فراگیری شیوه‌ی کنترل و مدیریت آن، خواهید دید که اعلان خط فرمان می‌تواند به ابزاری بسیار کاربردی تبدیل شود.

\section{ساختار اعلان خط فرمان (Prompt)}

اعلان پیش‌فرض در سیستم ما معمولاً به این صورت ظاهر می‌شود:

\begin{latin}
\begin{lstlisting}
[me@linuxbox ~]$
\end{lstlisting}
\end{latin}

دقت کنید که این رشته شامل نام کاربری، نام میزبان (\en{Hostname}) و دایرکتوری کاری فعلی (\en{Current Working Directory}) است؛ اما این ساختار چگونه شکل می‌گیرد؟ پاسخ در یک متغیر محیطی (\en{Environment Variable}) به نام \cmd{PS1} (مخفف \en{Prompt String 1}) نهفته است. می‌توانیم محتوای متغیر \cmd{PS1} را با استفاده از دستور \cmd{echo} مشاهده کنیم:

\begin{latin}
\begin{lstlisting}
[me@linuxbox ~]$ echo $PS1
[\u@\h \w]\$
\end{lstlisting}
\end{latin}

\begin{note}
\textbf{نکته:} در صورتی که خروجی سیستم شما با مثال فوق متفاوت بود، جای نگرانی نیست. هر توزیع لینوکس، رشته‌ی اعلان خط فرمان را بر اساس استانداردهای خود تعریف می‌کند و این تعاریف گاه تفاوت‌های چشمگیری با یکدیگر دارند.
\end{note}

خروجی حاصل نشان می‌دهد که متغیر \cmd{PS1} ترکیبی از کاراکترهای ثابت (مانند کروشه‌ها و نماد \cmd{@}) و کدهای کنترلی است که در نگاه نخست مبهم به نظر می‌رسند. این کدها در واقع «توالی‌های گریز» (\en{Escape Sequences}) هستند که با نویسه‌ی بک‌اسلش (\cmd{\textbackslash}) آغاز می‌شوند؛ مشابه آنچه در فصل هفتم بررسی کردیم. جدول زیر فهرستی از این کدهای ویژه را ارائه می‌دهد که \en{Bash} آن‌ها را در رشته‌ی اعلان خط فرمان به مقادیر پویا ترجمه می‌کند:

\begin{table}[htbp]
\centering
\caption{نویسه‌های کنترلی اعلان خط فرمان \en{Bash}}
\label{tab:bash-prompt-escapes}
\begin{tabularx}{\textwidth}{l X}
\toprule
\textbf{توالی گریز} & \textbf{مقدار جایگزین} \\
\midrule
\cmd{\textbackslash a} & نویسهٔ زنگ (\en{Bell}) یا هشدار صوتی. در گذشته برای به صدا درآوردن یک زنگ مکانیکی در تله‌تایپ‌ها استفاده می‌شد و امروزه معمولاً یک هشدار صوتی یا بصری (مانند چشمک زدن پنجره) را در ترمینال فعال می‌کند. \\
\addlinespace
\cmd{\textbackslash d} & تاریخ جاری با فرمت روز هفته، نام ماه و عدد روز (مثلاً \en{``Mon May 26''}). \\
\addlinespace
\cmd{\textbackslash h} & نام میزبان (\en{Hostname}) را تا اولین نقطه (\cmd{.}) نمایش می‌دهد. \\
\addlinespace
\cmd{\textbackslash H} & نام کامل میزبان (\en{Fully Qualified Domain Name}). \\
\addlinespace
\cmd{\textbackslash j} & تعداد وظایف (\en{Jobs}) فعال در نشست فعلی پوسته. \\
\addlinespace
\cmd{\textbackslash l} & نام پایه (\en{Basename}) دستگاه ترمینال جاری (مثلاً \en{``ttys0''}). \\
\addlinespace
\cmd{\textbackslash n} & نویسه‌ی خط جدید (\en{Newline}) برای ایجاد اعلان‌های چندخطی. \\
\addlinespace
\cmd{\textbackslash r} & بازگشت مکان‌نما به ابتدای سطر (\en{Carriage Return}). \\
\addlinespace
\cmd{\textbackslash s} & نام پوسته‌ای را که در حال اجراست، به شکلی که فراخوانی شده است، نمایش می‌دهد (مانند \cmd{bash} یا \cmd{sh}). \\
\addlinespace
\cmd{\textbackslash t} & زمان جاری با فرمت ۲۴ ساعته (\en{HH:MM:SS}). نمونه: \cmd{14:30:45} \\
\addlinespace
\cmd{\textbackslash T} & زمان جاری با فرمت ۱۲ ساعته (\en{HH:MM:SS}). نمونه: \cmd{02:30:45} \\
\addlinespace
\cmd{\textbackslash @} & زمان جاری با فرمت ۱۲ ساعته و همراه با \en{AM/PM}. نمونه: \cmd{02:30 PM} \\
\addlinespace
\cmd{\textbackslash A} & زمان جاری با فرمت ۲۴ ساعته (\en{HH:MM}). نمونه: \cmd{14:30} \\
\addlinespace
\cmd{\textbackslash u} & نام کاربری کاربر جاری که در سیستم احراز هویت کرده است را نمایش می‌دهد. \\
\addlinespace
\cmd{\textbackslash v} & شمارهٔ نسخهٔ (\en{version}) پوسته \en{Bash} را نمایش می‌دهد (مثلاً \cmd{2.00}). \\
\addlinespace
\cmd{\textbackslash V} & شمارهٔ انتشار (\en{release}) پوسته \en{Bash} را به همراه شمارهٔ سطح وصله (\en{patchlevel}) نمایش می‌دهد (مثلاً \cmd{2.00.0}). \\
\addlinespace
\cmd{\textbackslash w} & مسیر کامل دایرکتوری کاری فعلی، که در آن متغیر محیطی \cmd{\$HOME} با \cmd{\textasciitilde} جایگزین می‌شود. \\
\addlinespace
\cmd{\textbackslash W} & فقط نام آخرین دایرکتوری از مسیر کاری فعلی را نمایش می‌دهد و اگر در دایرکتوری خانگی (\cmd{\$HOME}) باشید، آن را با نویسه‌ی \cmd{\textasciitilde} خلاصه می‌کند. مثال: اگر مسیر جاری \file{home/rexbell/projects/} باشد، \cmd{\textbackslash W} مقدار \cmd{projects} را نشان می‌دهد. \\
\addlinespace
\cmd{\textbackslash !} & شمارهٔ دستور را در فهرست تاریخچهٔ کلی پوسته (\file{\textasciitilde/.bash\_history}) نشان می‌دهد. این شماره موقعیت دستور را در فایل تاریخچه (\en{History File}) مشخص می‌کند و پس از خروج و ورود مجدد، به‌صورت افزایشی و پیوسته باقی می‌ماند. \\
\addlinespace
\cmd{\textbackslash \#} & شمارهٔ دستور جاری را در نشست فعلی پوسته نمایش می‌دهد. این شماره موقعیت دستور را در نشست (\en{Session}) جاری پوسته مشخص می‌کند و از ابتدای هر نشست جدید از صفر شروع به شمارش می‌کند و فقط تا پایان همان نشست معتبر است. \\
\addlinespace
\cmd{\textbackslash \$} & اگر کاربر ریشه (\en{Root}) با شناسهٔ کاربری مؤثر (\en{Effective UID}) برابر صفر باشد، نویسهٔ \cmd{\#} را نمایش می‌دهد و در غیر این صورت، نویسهٔ \cmd{\$} را نمایش می‌دهد. \\
\addlinespace
\cmd{\textbackslash [} & شروع یک دنبالهٔ نویسه‌های غیرقابل چاپ (مانند کدهای کنترل ترمینال) را مشخص می‌کند. \\
\addlinespace
\cmd{\textbackslash ]} & پایان یک دنبالهٔ نویسه‌های غیرقابل چاپ را مشخص می‌کند. \\
\bottomrule
\end{tabularx}
\end{table}

\begin{note}
\textbf{نکته:} فرض کنید نام کامل میزبان (\en{FQDN}) شما \file{linuxbox.example.com} باشد. در این صورت، \cmd{\textbackslash h} فقط بخش \cmd{linuxbox} را نمایش می‌دهد (همان پیشوند یا نام کوتاه میزبان)، در حالی که \cmd{\textbackslash H} (حرف بزرگ) نام کامل یعنی \file{linuxbox.example.com} را نشان می‌دهد. استفاده از \cmd{\textbackslash h} باعث می‌شود اعلان خط فرمان کوتاه‌تر و خوانا‌تر باقی بماند.
\end{note}

\section{آزمودن الگوهای اعلان خط فرمان}

با بهره‌گیری از فهرست نویسه‌های کنترلی، می‌توانیم ساختار اعلان خط فرمان (\en{Prompt}) را تغییر داده و تأثیر آن را بلافاصله مشاهده کنیم. در گام نخست، از رشته‌ی فعلی متغیر \cmd{PS1} نسخه‌ی پشتیبان تهیه می‌کنیم تا در صورت نیاز، امکان بازیابی آن فراهم باشد. برای این منظور، مقدار فعلی را در یک متغیر شخصی ذخیره می‌کنیم:

\begin{latin}
\begin{lstlisting}
[me@linuxbox ~]$ ps1_old="$PS1"
\end{lstlisting}
\end{latin}

در اینجا متغیری به نام \cmd{ps1\_old} ایجاد کرده و مقدار \cmd{PS1} را به آن تخصیص داده‌ایم. با استفاده از دستور \cmd{echo} می‌توانیم از صحت کپی شدن رشته اطمینان حاصل کنیم:

\begin{latin}
\begin{lstlisting}
[me@linuxbox ~]$ echo $ps1_old
[\u@\h \w]\$
\end{lstlisting}
\end{latin}

هر زمان که در طول این نشستِ ترمینال بخواهیم، می‌توانیم با معکوس کردن فرآیند فوق، اعلان خط فرمان اصلی را بازیابی کنیم:

\begin{latin}
\begin{lstlisting}
[me@linuxbox ~]$ PS1="$ps1_old"
\end{lstlisting}
\end{latin}

اکنون برای شروع آزمایش، بیایید بررسی کنیم که در صورت تهی کردن رشته‌ی اعلان، چه اتفاقی رخ می‌دهد:

\begin{latin}
\begin{lstlisting}
[me@linuxbox ~]$ PS1=
\end{lstlisting}
\end{latin}

اگر هیچ مقداری به متغیر \cmd{PS1} اختصاص ندهیم، اعلان خط فرمان ناپدید می‌شود! در این حالت پوسته همچنان فعال و آماده‌ی دریافت فرامین است، اما هیچ نشانه‌ی بصری نمایش نمی‌دهد. از آنجا که کار در این وضعیت دشوار است، آن را با یک اعلان حداقلی جایگزین می‌کنیم:

\begin{latin}
\begin{lstlisting}
PS1="\$ "
\end{lstlisting}
\end{latin}

اکنون وضعیت بهبود یافت و دست‌کم نقطه‌ی درج دستورات قابل رویت است. به فاصله‌ی (\en{Space}) انتهایی داخل علامت‌های نقل قول دقت کنید؛ این فاصله باعث ایجاد یک فضای خالی میان علامت دلار و مکان‌نما (\en{Cursor}) می‌شود تا خوانایی دستورات افزایش یابد.

بیایید یک هشدار صوتی به اعلان خط فرمان اضافه کنیم:

\begin{latin}
\begin{lstlisting}
$ PS1="[\a]\$ "
\end{lstlisting}
\end{latin}

با این تنظیم، هر بار که اعلان خط فرمان ظاهر شود، باید صدای بوق (\en{Beep}) شنیده شود؛ هرچند در برخی توزیع‌ها این قابلیت به‌صورت پیش‌فرض غیرفعال شده است. اگرچه ممکن است این وضعیت در درازمدت آزاردهنده باشد، اما برای آگاهی از اتمام اجرای دستورات طولانی بسیار کاربردی است. دقت کنید که از توالی‌های \cmd{\textbackslash [} و \cmd{\textbackslash ]} استفاده کرده‌ایم؛ از آنجا که کاراکتر زنگ (\cmd{\textbackslash a}) یک نویسه غیرقابل چاپ است و تأثیری در جابه‌جایی مکان‌نما ندارد، باید با این علائم به \en{Bash} اطلاع دهیم تا طول رشته‌ی اعلان را برای مدیریت صحیح فضای سطر، درست محاسبه کند.

حالا یک اعلان جدید طراحی می‌کنیم که شامل زمان جاری و نام میزبان باشد:

\begin{latin}
\begin{lstlisting}
$ PS1="\A \h \$ "
17:33 linuxbox $
\end{lstlisting}
\end{latin}

درج زمان در اعلان خط فرمان برای پایش مدت‌زمان اجرای وظایف یا ثبت زمان دقیق فعالیت‌ها بسیار مفید است. در نهایت، اعلانی مشابه ساختار اصلی، اما با ظاهری متفاوت ایجاد می‌کنیم:

\begin{latin}
\begin{lstlisting}
17:37 linuxbox $ PS1="<\u@\h \w>\$ "
<me@linuxbox ->$
\end{lstlisting}
\end{latin}

پیشنهاد می‌شود سایر توالی‌های معرفی‌شده در جدول قبلی را نیز آزمایش کنید تا به ترکیبی خلاقانه و متناسب با نیازهای کاری خود دست یابید.

\section{اعمال رنگ‌بندی در اعلان خط فرمان}

اکثریت نرم‌افزارهای شبیه‌ساز ترمینال (\en{Terminal Emulators}) قابلیت تفسیر توالی‌های گریز (\en{escape sequences}) غیرقابل چاپ را دارند. این توالی‌ها به‌منظور مدیریت ویژگی‌های بصری نویسه‌ها (نظیر رنگ، ضخامت متن، و حالت چشمک‌زن) و همچنین کنترل دقیق موقعیت مکان‌نما (\en{Cursor}) به‌کار می‌روند. در ادامه، ابتدا به چگونگی اعمال رنگ‌ها خواهیم پرداخت و سپس نحوه‌ی مدیریت موقعیت مکان‌نما را بررسی می‌کنیم.

\subsection{تکامل ترمینال‌ها و استانداردهای کنترلی}

در دوران پیشین، زمانی که ترمینال‌ها (\en{Terminals}) به رایانه‌های میزبانِ از راه‌دور متصل می‌شدند، تنوع گسترده‌ای از برندهای تجاری وجود داشت که هر یک با پروتکل‌های عملیاتی متفاوتی کار می‌کردند. این دستگاه‌ها نه تنها در چیدمان صفحه‌کلید، بلکه در شیوه‌ی تفسیر داده‌های کنترلی نیز با یکدیگر تفاوت داشتند. برای مدیریت این ناهمگونی‌ها، سیستم‌عامل‌های خانواده یونیکس دو زیرسیستم (\en{subsystem}) پیچیده به نام‌های \cmd{termcap} و \cmd{terminfo} را توسعه دادند. امروزه نیز اگر تنظیمات نرم‌افزار شبیه‌ساز ترمینال خود را بررسی کنید، احتمالاً با گزینه‌هایی برای تعیین نوع شبیه‌سازی (مانند \cmd{xterm} یا \cmd{vt100}) مواجه خواهید شد.

به منظور ایجاد یک زبان مشترک میان ترمینال‌ها، مؤسسه ملی استاندارد آمریکا (\en{ANSI}) مجموعه‌ای استاندارد از توالی‌های گریز (\en{escape sequences}) را برای کنترل نمایشگرهای متنی تدوین کرد. کاربران باسابقه سیستم‌عامل \en{DOS} احتمالاً فایل \file{ANSI.SYS} را به‌خاطر دارند که وظیفه‌ی فعال‌سازی و تفسیر این کدهای استاندارد را بر عهده داشت.

رنگ متن از طریق ارسال «توالی‌های گریز \en{ANSI}» به شبیه‌ساز ترمینال کنترل می‌شود. این کدها در میان جریان نویسه‌های متنی جاسازی شده و مستقیماً روی نمایشگر چاپ نمی‌شوند؛ بلکه ترمینال آن‌ها را به عنوان دستورالعمل‌های قالب‌بندی تفسیر می‌کند. همان‌طور که پیش‌تر ذکر شد، از توالی‌های \cmd{\textbackslash [} و \cmd{\textbackslash ]} برای محصور کردن نویسه‌های غیرقابل چاپ استفاده می‌کنیم تا اختلالی در محاسبه‌ی طول سطر ایجاد نشود.

یک کد گریز \en{ANSI} با مقدار هشت‌تایی (\en{Octal}) \cmd{033} آغاز می‌شود (کدی که با فشردن کلید \en{Esc} تولید می‌شود)، و پس از آن یک ویژگی نویسه (\en{Character Attribute}) به‌صورت اختیاری می‌آید و در ادامه یک دستور قرار می‌گیرد.

برای مثال، کد تنظیم رنگ متن به حالت عادی (ویژگی = 0) و رنگ مشکی، به این صورت است:

\begin{latin}
\begin{lstlisting}
\033[0;30m
\end{lstlisting}
\end{latin}

جدول زیر فهرست رنگ‌های متنی موجود را نشان می‌دهد. توجه کنید که رنگ‌ها به دو گروه تقسیم می‌شوند که با اعمال ویژگی پررنگ یا \en{Bold} (با کد 1) از هم متمایز می‌شوند؛ این ویژگی است که ظاهر رنگ‌های «روشن» (\en{Light}) را ایجاد می‌کند.

\begin{table}[htbp]
\centering
\caption{توالی‌های گریز \en{ANSI} برای تنظیم رنگ متن}
\label{tab:ansi-fg}
\begin{tabularx}{\textwidth}{l l l l}
\toprule
\textbf{توالی گریز (رنگ عادی)} & \textbf{رنگ متن} & \textbf{توالی گریز (رنگ روشن)} & \textbf{رنگ متن} \\
\midrule
\cmd{\textbackslash 033[0;30m} & مشکی & \cmd{\textbackslash 033[1;30m} & خاکستری تیره \\
\addlinespace
\cmd{\textbackslash 033[0;31m} & قرمز & \cmd{\textbackslash 033[1;31m} & قرمز روشن \\
\addlinespace
\cmd{\textbackslash 033[0;32m} & سبز & \cmd{\textbackslash 033[1;32m} & سبز روشن \\
\addlinespace
\cmd{\textbackslash 033[0;33m} & قهوه‌ای & \cmd{\textbackslash 033[1;33m} & زرد \\
\addlinespace
\cmd{\textbackslash 033[0;34m} & آبی & \cmd{\textbackslash 033[1;34m} & آبی روشن \\
\addlinespace
\cmd{\textbackslash 033[0;35m} & بنفش & \cmd{\textbackslash 033[1;35m} & بنفش روشن \\
\addlinespace
\cmd{\textbackslash 033[0;36m} & فیروزه‌ای & \cmd{\textbackslash 033[1;36m} & فیروزه‌ای روشن \\
\addlinespace
\cmd{\textbackslash 033[0;37m} & خاکستری روشن & \cmd{\textbackslash 033[1;37m} & سفید \\
\bottomrule
\end{tabularx}
\end{table}

بیایید یک اعلان خط فرمان با رنگ‌بندی قرمز طراحی کنیم. برای این منظور، توالی گریز (\en{Escape Sequence}) مربوط به رنگ را در ابتدای رشته‌ی متغیر \cmd{PS1} قرار می‌دهیم:

\begin{latin}
\begin{lstlisting}
<me@linuxbox ->$ PS1="\[\033[0;31m\]<\u@\h \W>\$ "
<me@linuxbox ->$
\end{lstlisting}
\end{latin}

همان‌طور که مشاهده می‌کنید، تغییر رنگ با موفقیت اعمال شد؛ اما یک ایراد فنی وجود دارد: تمامی دستوراتی که پس از اعلان تایپ می‌کنیم نیز به رنگ قرمز نمایش داده می‌شوند. برای حل این مسئله، باید یک کد گریز دیگر را در انتهای رشته‌ی اعلان اضافه کنیم تا به شبیه‌ساز ترمینال فرمان بازگشت به حالت پیش‌فرض (\en{Reset}) را صادر کند:

\begin{latin}
\begin{lstlisting}
<me@linuxbox ->$ PS1="\[\033[0;31m\]<\u@\h \W>\$\[\033[0m\] "
<me@linuxbox ->$
\end{lstlisting}
\end{latin}

اکنون خروجی کاملاً بهینه است؛ اعلان با رنگ قرمز متمایز شده و متن ورودی کاربر به حالت استاندارد بازگشته است.

علاوه بر رنگ قلم، امکان تنظیم رنگ پس‌زمینه (\en{Background}) برای متن نیز با استفاده از کدهای مندرج در جدول زیر فراهم است. شایان ذکر است که رنگ‌های پس‌زمینه برخلاف رنگ متن، از ویژگی پررنگ (\en{Bold}) پشتیبانی نمی‌کنند.

\begin{table}[htbp]
\centering
\caption{کدهای \en{ANSI} برای رنگ پس‌زمینه متن}
\label{tab:ansi-bg}
\begin{tabularx}{\textwidth}{l X}
\toprule
\textbf{توالی گریز} & \textbf{رنگ پس‌زمینه متن} \\
\midrule
\cmd{\textbackslash 033[0;40m} & مشکی \\
\addlinespace
\cmd{\textbackslash 033[0;41m} & قرمز \\
\addlinespace
\cmd{\textbackslash 033[0;42m} & سبز \\
\addlinespace
\cmd{\textbackslash 033[0;43m} & قهوه‌ای \\
\addlinespace
\cmd{\textbackslash 033[0;44m} & آبی \\
\addlinespace
\cmd{\textbackslash 033[0;45m} & بنفش \\
\addlinespace
\cmd{\textbackslash 033[0;46m} & فیروزه‌ای \\
\addlinespace
\cmd{\textbackslash 033[0;47m} & خاکستری روشن \\
\bottomrule
\end{tabularx}
\end{table}

با اعمال تغییری کوچک در نخستین توالی گریز، می‌توانیم اعلانی با پس‌زمینه قرمز ایجاد کنیم:

\begin{latin}
\begin{lstlisting}
<me@linuxbox ->$ PS1="\[\033[0;41m\]<\u@\h \W>\$\[\033[0m\] "
<me@linuxbox ->$
\end{lstlisting}
\end{latin}

پیشنهاد می‌شود کدهای مختلف را ترکیب کنید تا به طراحی مطلوب خود دست یابید.

\begin{note}
\textbf{نکته:} افزون بر ویژگی‌های «عادی» (0) و «پررنگ» (1)، می‌توانید از ویژگی‌های دیگری نظیر «زیرخط» (4)، «چشمک‌زن» (5) و «رنگ معکوس» (7) نیز استفاده کنید. با این حال، به دلیل ملاحظات بصری و تجربه کاربری، بسیاری از شبیه‌سازهای مدرن ترمینال از ویژگی چشمک‌زن (\en{Blink}) پشتیبانی نمی‌کنند.
\end{note}

\section{کنترل موقعیت مکان‌نما در ترمینال}

توالی‌های گریز (\en{Escape Sequences}) علاوه بر مدیریت رنگ، برای تعیین دقیق موقعیت مکان‌نما (\en{Cursor}) نیز به کار می‌روند. این قابلیت به‌طور معمول برای نمایش اطلاعات پویا، مانند ساعت یا اطلاعات لحظه‌ای سیستم، در نقاط خاصی از نمایشگر (مثلاً گوشه‌ی بالایی سمت راست) هم‌زمان با بازترسیم اعلان خط فرمان استفاده می‌شود. جدول زیر کدهای کاربردی جهت هدایت مکان‌نما را نمایش می‌دهد:

\begin{table}[htbp]
\centering
\caption{توالی‌های کنترلی حرکت مکان‌نما}
\label{tab:cursor-control}
\begin{tabularx}{\textwidth}{l X}
\toprule
\textbf{توالی گریز} & \textbf{عملکرد} \\
\midrule
\cmd{\textbackslash 033[L;CH} & مکان‌نما را به سطر \en{L} و ستون \en{C} منتقل می‌کند. \\
\addlinespace
\cmd{\textbackslash 033[nA} & مکان‌نما را به‌اندازهٔ \en{n} سطر به سمت بالا جابه‌جا می‌کند. \\
\addlinespace
\cmd{\textbackslash 033[nB} & مکان‌نما را به‌اندازهٔ \en{n} سطر به سمت پایین جابه‌جا می‌کند. \\
\addlinespace
\cmd{\textbackslash 033[nC} & مکان‌نما را به‌اندازهٔ \en{n} نویسه به جلو جابه‌جا می‌کند. \\
\addlinespace
\cmd{\textbackslash 033[nD} & مکان‌نما را به‌اندازهٔ \en{n} نویسه به عقب جابه‌جا می‌کند. \\
\addlinespace
\cmd{\textbackslash 033[2J} & صفحه را به‌طور کامل پاک می‌کند و مکان‌نما را به مبدأ (سطر 0، ستون 0) بازمی‌گرداند. \\
\addlinespace
\cmd{\textbackslash 033[K} & محتوای خط را از موقعیت فعلی مکان‌نما تا انتهای سطر پاک می‌کند. \\
\addlinespace
\cmd{\textbackslash 033[s} & مختصات فعلی مکان‌نما را در حافظه ذخیره می‌کند. \\
\addlinespace
\cmd{\textbackslash 033[u} & مکان‌نما را به مختصات ذخیره‌شده بازمی‌گرداند. \\
\bottomrule
\end{tabularx}
\end{table}

با ترکیب کدهای مندرج در جدول فوق، می‌توانیم اعلانی طراحی کنیم که در هر بار فراخوانی، یک نوار وضعیت قرمز رنگ در بالاترین بخش ترمینال ترسیم کرده و ساعت سیستم را با رنگ زرد در آن نمایش دهد. ساختار این رشته‌ی کنترلی به شرح زیر است:

\begin{latin}
\begin{lstlisting}
PS1="\[\033[s\033[0;0H\033[0;41m\033[K\033[1;33m\t\033[0m\033[u\]<\u@\h \W>\$ "
\end{lstlisting}
\end{latin}

جدول زیر به‌صورت دقیق وظیفه و عملکرد هر مؤلفه در این رشته‌ی پیچیده را تشریح می‌کند:

\begin{table}[htbp]
\centering
\caption{کالبدشکافی اجزای اعلان خط فرمان ترکیبی}
\label{tab:prompt-anatomy}
\begin{tabularx}{\textwidth}{l X}
\toprule
\textbf{توالی} & \textbf{عملکرد فنی و کاربرد} \\
\midrule
\cmd{\textbackslash [} & آغاز بخش نویسه‌های غیرقابل چاپ را نشان می‌دهد تا \en{Bash} طول واقعی این نویسه‌ها را در محاسبهٔ طول اعلان لحاظ نکند. \\
\addlinespace
\cmd{\textbackslash 033[s} & مختصات فعلی مکان‌نما (موقعیت آن در خط فرمان) را پیش از هر جابه‌جایی، در حافظه ذخیره می‌کند. \\
\addlinespace
\cmd{\textbackslash 033[0;0H} & مکان‌نما را بلافاصله به گوشهٔ بالا-چپ صفحهٔ نمایش (سطر اول، ستون اول) منتقل می‌کند. \\
\addlinespace
\cmd{\textbackslash 033[0;41m} & رنگ پس‌زمینه را قرمز تنظیم می‌کند (کد 41 رنگ پس‌زمینه، کد 0 بازنشانی سایر ویژگی‌ها). \\
\addlinespace
\cmd{\textbackslash 033[K} & محتوای سطر را از موقعیت فعلی مکان‌نما تا انتهای خط پاک می‌کند؛ از آنجا که پس‌زمینه از قبل قرمز شده، این عمل عملاً یک نوار قرمزرنگ در سراسر سطر بالا رسم می‌کند. \\
\addlinespace
\cmd{\textbackslash 033[1;33m} & رنگ متن (پیش‌زمینه) را به زرد پررنگ تغییر می‌دهد تا محتوای بعدی با این رنگ نمایش داده شود. \\
\addlinespace
\cmd{\textbackslash t} & ساعت جاری سیستم را نمایش می‌دهد. از آنجا که این نویسه در بالای صفحه چاپ می‌شود و سپس مکان‌نما به موقعیت اصلی بازمی‌گردد، طول آن در محاسبهٔ عرض اعلان اصلی خط فرمان دخالتی ندارد. \\
\addlinespace
\cmd{\textbackslash 033[0m} & تمام تنظیمات رنگی و ویژگی‌های قبلی را به حالت پیش‌فرض ترمینال بازمی‌گرداند. \\
\addlinespace
\cmd{\textbackslash 033[u} & مکان‌نما را به همان مختصاتی که پیش‌تر ذخیره شده بود، بازمی‌گرداند تا کاربر بتواند دستور خود را در محل درست وارد کند. \\
\addlinespace
\cmd{\textbackslash ]} & پایان بخش نویسه‌های غیرقابل چاپ را اعلام می‌کند. \\
\addlinespace
\cmd{<\textbackslash u@\textbackslash h \textbackslash W>\$} & بدنهٔ اصلی و قابل‌مشاهدهٔ اعلان است که نام کاربری، نام میزبان، و نام دایرکتوری کاری جاری را نمایش می‌دهد. \\
\bottomrule
\end{tabularx}
\end{table}

\section{ذخیره‌سازی پیکربندی اعلان خط فرمان}

بدیهی است که تایپ مکرر چنین رشته‌های پیچیده‌ای منطقی نیست؛ بنابراین برای استفاده‌ی همیشگی، باید آن را در فایل تنظیمات پوسته ذخیره کنیم. برای دائمی کردن اعلان اختصاصی خود، کافی است خطوط زیر را به انتهای فایل \file{.bashrc} در دایرکتوری خانگی (\en{Home}) اضافه کنید:

\begin{latin}
\begin{lstlisting}
PS1="\[\033[s\033[0;0H\033[0;41m\033[K\033[1;33m\t\033[0m\033[u\]<\u@\h \W>\$ "

export PS1
\end{lstlisting}
\end{latin}

\section{جمع‌بندی}

شاید باورکردنی نباشد، اما قابلیت‌های اعلان خط فرمان فراتر از این موارد است و با استفاده از توابع (\en{Functions}) و اسکریپت‌های پوسته، می‌توان اعلان‌هایی هوشمندتر و پویاتر طراحی کرد. با این حال، مفاهیم مطرح شده در این فصل نقطه‌ی آغاز ایده‌آلی محسوب می‌شوند. اگرچه بسیاری از کاربران به دلیل کارآمدی اعلان پیش‌فرض تمایلی به تغییر آن ندارند، اما برای علاقه‌مندان به شخصی‌سازی (\en{Customization})، این مبحث پتانسیل بالایی برای کاوش و خلاقیت فراهم می‌کند.

\section{منابع مطالعه تکمیلی}

مستند \en{``The Bash Prompt HOWTO''} از پروژه مستندات لینوکس (\en{LDP})، مبحث قابلیت‌های اعلان خط فرمان (\en{shell prompt}) را به‌طور جامع پوشش می‌دهد:

\begin{latin}
\url{http://tldp.org/HOWTO/Bash-Prompt-HOWTO/}
\end{latin}

مقاله ویکی‌پدیا درباره کدهای گریز \en{ANSI} (\en{ANSI Escape Codes}):

\begin{latin}
\url{http://en.wikipedia.org/wiki/ANSI_escape_code}
\end{latin}